\chapter{Loans, CLOs and Securitisation}\label{ch:m2:loans-clos-and-securitisation}

At the end of 2025 American companies owed USD~1.55 trillion of institutional \omterm{def:m2:loans-clos-and-securitisation:loan}{leveraged
loans}, 9.2\% more than a year before, by the Federal Reserve's count. These are loans
to companies with high debt, and they are not kept by the banks that arrange them:
they are sold to institutional investors, and many end up in vehicles that buy a
diversified pool of them and pay for it by selling notes of different seniority. Each vehicle's documents say, to the dollar, who is paid first each
quarter, and contain tests that, when the loans start to default, turn off the cash to
the most junior investors and send it to the most senior. This chapter explains the
loans, \omterm{def:m2:loans-clos-and-securitisation:sec}{securitisation}, the \omterm{def:m2:loans-clos-and-securitisation:clo}{collateralised loan obligation} and its waterfall, and who
holds which part; the build is the waterfall itself.

\section{The leveraged-loan market}

\begin{definition}[Leveraged loan, term loan B, covenant-lite]\label{def:m2:loans-clos-and-securitisation:loan}
A \emph{leveraged loan}\index{leveraged loan} is a loan, usually senior and secured on
the borrower's assets and paying a floating rate plus a spread, to a company whose
debt is high relative to its earnings or whose credit is rated below investment grade.
A \emph{term loan B}\index{term loan B} is the tranche of such a loan arranged to be
sold to institutional investors: a long maturity, little amortisation before it, and
tradable. A \emph{covenant-lite}\index{covenant-lite} loan has no maintenance
\omterm{def:m2:corporate-bonds:covenant}{covenants}, financial ratios the borrower must meet at every test date; it keeps only
incurrence \omterm{def:m2:corporate-bonds:covenant}{covenants}, which are tested when the borrower takes an action such as
borrowing more.
\end{definition}

A bank arranges the loan, for example to finance a private-equity firm's purchase of
a company, and syndicates it: it sells most of it to loan funds, CLOs and other
investors, who then trade it in a secondary market. Because the loan floats, the
borrower's interest bill moves with the policy rate of chapter 1; because the loan is
senior and usually secured, its lenders tend to recover more after a default than
bondholders do (chapter 23's auctions show it); and when it is \omterm{def:m2:loans-clos-and-securitisation:loan}{covenant-lite}, they
have fewer early warnings before one.

The Federal Reserve watches this market in its financial stability reports. In May
2026 it noted that the share of new loans to companies with debt of four or more times
earnings had risen and stayed above its historical median, that the median borrower's
interest coverage was near its historical low, and that defaults including distressed
exchanges, loans renegotiated when the borrower was in difficulty, remained relatively
high.

\begin{dated}{2026-09}{The leveraged-loan and CLO markets}\label{dat:m2:loans-clos-and-securitisation:size}
Institutional \omterm{def:m2:loans-clos-and-securitisation:loan}{leveraged loans} outstanding in the United States: USD~1\,549 billion at
the end of 2025, up 9.2\% in a year and by 12.6\% a year on average since 2000
(Federal Reserve, Financial Stability Report, May 2026, from PitchBook LCD data).
CLOs: USD~264 billion in 2011, USD~617 billion in 2018 (Federal Reserve, 2019, from
SIFMA data); no later public figure was fetched for this book.
\end{dated}

\section{Securitisation}

\begin{definition}[Securitisation, special-purpose vehicle]\label{def:m2:loans-clos-and-securitisation:sec}
\emph{Securitisation}\index{securitisation} is the pooling of loans or other
financial assets and the sale of claims on their cash flows as securities. A
\emph{special-purpose vehicle}\index{special-purpose vehicle} (SPV) is a company set up
only to hold the pool and issue the securities, so that the pool is separate from
whoever originated or sold the assets and the investors' claims depend only on it.
\end{definition}

\omterm{def:m2:loans-clos-and-securitisation:sec}{Securitisation} turns a pool of assets into securities with different risks. The
mortgage pass-through of chapter 12 passes the cash flows through pro rata; most other
\omterm{def:m2:loans-clos-and-securitisation:sec}{securitisations} tranche them, as the index tranches of chapter 24 do, so that some
investors are paid first and others absorb the first losses. Credit cards, car loans,
commercial mortgages and \omterm{def:m2:loans-clos-and-securitisation:loan}{leveraged loans} are all securitised. The Federal Reserve
notes that \omterm{def:m2:loans-clos-and-securitisation:sec}{securitisation} can add leverage to the financial system, because the
vehicles are subject to lighter rules, such as risk retention, than banks' capital
requirements.

Risk retention requires the sponsor of a \omterm{def:m2:loans-clos-and-securitisation:sec}{securitisation} to keep a share of its credit
risk, 5\% under the American rule adopted under the Dodd--Frank Act, so that it has
a stake in the quality of what it sells. For one kind of CLO, it does not apply: in February 2018 the
federal appeals court in Washington held that managers of open-market CLOs, who
neither originate the loans nor hold them before the vehicle does, are not the
securitisers the rule covers, and vacated the rule as applied to them.

\section{The CLO waterfall and its tests}

\begin{definition}[Collateralised loan obligation, payment waterfall]\label{def:m2:loans-clos-and-securitisation:clo}
A \emph{collateralised loan obligation}\index{collateralised loan obligation} (CLO)
is an SPV that holds an actively managed, diversified portfolio of \omterm{def:m2:loans-clos-and-securitisation:loan}{leveraged loans}
and funds it by issuing rated notes of decreasing seniority and unrated equity. Its
\emph{payment waterfall}\index{payment waterfall} is the order, fixed in its documents,
in which the cash collected from the loans is paid out on each payment date.
\end{definition}

\begin{definition}[Overcollateralisation test, interest-coverage test, equity tranche]\label{def:m2:loans-clos-and-securitisation:tests}
An \emph{overcollateralisation test}\index{overcollateralisation test} compares the
par value of the loans with the notes outstanding down to a given class; an
\emph{interest-coverage test}\index{interest-coverage test} compares the interest
collected with the interest due on those notes. Each has a trigger ratio, and when a
ratio falls below it, cash that would have gone further down the waterfall repays the
most senior notes instead. The \emph{equity tranche}\index{equity tranche} of a
\omterm{def:m2:loans-clos-and-securitisation:sec}{securitisation} is its most junior claim, unrated, which receives whatever cash is left
after every other claim has been paid.
\end{definition}

A manager selects the loans, trades them, and during a reinvestment period of several
years replaces those that repay or default; the notes pay a floating rate plus a
spread fixed at issue, and the equity is paid the difference between what the loans
earn and what the notes cost, the \emph{excess spread}, as long as the tests pass
(\cref{fig:m2:loans-clos-and-securitisation:structure}).

\begin{omfigure}
\centering
\begin{tikzpicture}[font=\small,
  box/.style={draw,thick,rounded corners=2pt,align=center,minimum height=0.9cm,inner sep=3pt},
  arr/.style={-{Stealth[length=2.5mm]},thick}]
\node[box,draw=gray,fill=gray!8,minimum width=3cm] (loans) at (0,0) {a diversified pool\\of leveraged loans\\(USD~500 million par)};
\node[box,draw=omThm,fill=omThm!8,minimum width=2.6cm] (spv) at (5.4,0) {CLO (SPV)\\run by a manager};
\draw[arr] (loans) -- node[above,font=\footnotesize] {interest} node[below,font=\footnotesize] {and principal} (spv);
\foreach \n/\t/\y/\c in {aaa/{AAA 310 (62\%)}/1.6/omDef, aa/{AA 55 (11\%)}/0.8/omDef, a/{A 30 (6\%)}/0/omProp, bbb/{BBB 30 (6\%)}/-0.8/omProp, bb/{BB 25 (5\%)}/-1.6/omProp}
  \node[box,draw=\c,fill=\c!8,minimum width=3.2cm,minimum height=0.6cm] (\n) at (10.2,\y) {\t};
\node[box,draw=gray,fill=gray!25,minimum width=3.2cm,minimum height=0.6cm] (eq) at (10.2,-2.4) {equity 50 (10\%)};
\foreach \n in {aaa,aa,a,bbb,bb,eq} \draw[arr] (spv.east) -- (\n.west);
\node[font=\footnotesize,align=left,anchor=west] at (12.0,1.2) {paid first,\\lowest spread};
\node[font=\footnotesize,align=left,anchor=west] at (12.0,-2.0) {paid last,\\first loss};
\end{tikzpicture}

\omcaption{The illustrative CLO of this chapter: a managed pool of \omterm{def:m2:loans-clos-and-securitisation:loan}{leveraged loans} held
by a \omterm{def:m2:loans-clos-and-securitisation:sec}{special-purpose vehicle} and funded by five classes of rated notes and 10\%
equity, in USD millions. Cash is paid down the stack in order; losses are taken up
from the bottom. Schematic; the structure is illustrative.}\label{fig:m2:loans-clos-and-securitisation:structure}
\end{omfigure}

Each quarter the interest collected pays the manager's senior fee, then the interest
on each class in turn. After a class is paid, its tests are checked. If one fails, the
interest left is used to repay the senior notes until the ratio is back at its trigger,
and only what then remains continues down
(\cref{fig:m2:loans-clos-and-securitisation:waterfall}). Unpaid interest on the junior
notes is not a default: it is added to their balance and paid later if it can be.

\begin{omfigure}
\centering
\begin{tikzpicture}[font=\small,
  wstep/.style={draw,thick,rounded corners=2pt,align=center,minimum width=4.4cm,minimum height=0.62cm,inner sep=2pt},
  test/.style={draw,thick,diamond,aspect=2.6,align=center,inner sep=1pt,font=\footnotesize},
  arr/.style={-{Stealth[length=2.2mm]},thick}]
\node[wstep,draw=gray,fill=gray!8] (in) at (0,0) {interest collected};
\node[wstep,draw=gray,fill=gray!8] (fee) at (0,-1.0) {senior fee};
\node[wstep,draw=omDef,fill=omDef!8] (s) at (0,-2.0) {AAA, then AA interest};
\node[test,draw=omThm] (t1) at (0,-3.25) {AA tests\\pass?};
\node[wstep,draw=omProp,fill=omProp!8] (m) at (0,-4.5) {A, BBB, BB interest,\\each followed by its test};
\node[wstep,draw=gray,fill=gray!25] (eq) at (0,-5.7) {equity: what is left};
\node[wstep,draw=omThm,fill=omThm!8,minimum width=3.6cm] (div) at (6,-3.25) {repay AAA notes until\\the test is cured};
\draw[arr] (in) -- (fee); \draw[arr] (fee) -- (s); \draw[arr] (s) -- (t1);
\draw[arr] (t1) -- node[left,font=\footnotesize] {yes} (m); \draw[arr] (m) -- (eq);
\draw[arr] (t1) -- node[above,font=\footnotesize] {no} (div);
\draw[arr] (div.south) |- node[pos=0.25,right,font=\footnotesize,align=left] {rest, if any,\\continues} (m.east);
\draw[arr,dashed,omThm] (m.west) -- ++(-1.4,0) |- node[pos=0.25,left,font=\footnotesize,align=right] {junior test\\fails: divert} (s.west);
\end{tikzpicture}

\omcaption{The interest waterfall of the chapter's CLO. After each class is paid, its
tests are checked; a failure diverts the remaining interest to repay the senior notes
until the test is cured, and only what is left continues down to the equity.
Schematic.}\label{fig:m2:loans-clos-and-securitisation:waterfall}
\end{omfigure}

\begin{example}[The excess spread]\label{ex:m2:loans-clos-and-securitisation:base}
The chapter's CLO holds USD~500 million of loans paying 3.5\% over a floating rate of
4\%, USD~37.5 million a year. The manager's fee of 0.45\% costs USD~2.25 million, and
the notes, from AAA at 1.30\% over to BB at 6.00\% over, cost USD~26.14 million. The
equity, USD~50 million, receives USD~9.11 million a year, 18.2\% of its investment;
without defaults and with the loans sold at par after five years, its internal rate of
return is 19.5\%. The overcollateralisation ratios start at 1.370, 1.266, 1.176 and
1.111 at the AA, A, BBB and BB levels, against triggers of 1.25, 1.18, 1.10 and 1.05.
\end{example}

\begin{proposition}[The overcollateralisation cushion]\label{prop:m2:loans-clos-and-securitisation:cushion}
If the notes down to a class total $D$, the loans have par $P$ and the class's trigger
is $\tau$, the test fails once the par has fallen by more than $P - \tau D$. With
recoveries reinvested at par, each unit of defaulted par lowers the portfolio's par by
$1 - R$, so the test fails once cumulative defaults exceed $(P - \tau D)/(1 - R)$.
\end{proposition}

\begin{proof}
The test is $P'/D \geq \tau$. Before any diversion $D$ is unchanged, so it fails when
$P' < \tau D$, a fall of $P - \tau D$. A default of $x$ removes $x$ of par and the
reinvested recovery adds back $Rx$, a net fall of $(1-R)x$.
\end{proof}

For the BB test, $P - \tau D = 500 - 1.05 \times 450 = 27.5$: with a recovery of 70\%,
USD~91.7 million of defaults, 18.3\% of the portfolio, trip it. Once it fails, the cure
takes roughly the par lost each quarter, divided by the trigger, out of the interest
that would have gone down the waterfall; when that exceeds what is left after the notes'
interest, the equity receives nothing
(\cref{fig:m2:loans-clos-and-securitisation:paths}).

\begin{omfigure}
\begin{tikzpicture}
\begin{axis}[width=11cm,height=4.6cm,
  xlabel={quarter},ylabel={BB-level OC ratio},
  tick label style={font=\small},label style={font=\small},
  xmin=0.5,xmax=20.5,ymin=1.0,ymax=1.12,grid=major,grid style={gray!25},
  yticklabel style={/pgf/number format/.cd,fixed,precision=2},
  legend columns=4,legend style={font=\footnotesize,at={(0.5,-0.32)},anchor=north,draw=none,fill=none,/tikz/every even column/.append style={column sep=8pt}},
  table/col sep=comma]
\addplot[omDef,very thick] table[x=quarter,y=oc0]{figdata/markets-2/25-loans-clos-and-securitisation/paths.csv};
\addplot[omThm,very thick] table[x=quarter,y=oc1]{figdata/markets-2/25-loans-clos-and-securitisation/paths.csv};
\addplot[omProp,very thick] table[x=quarter,y=oc2]{figdata/markets-2/25-loans-clos-and-securitisation/paths.csv};
\addplot[black,thick,dashed] coordinates {(0.5,1.05) (20.5,1.05)};
\legend{2\% a year,4.8\%,7\%,trigger 1.05}
\end{axis}
\end{tikzpicture}

\begin{tikzpicture}
\begin{axis}[width=11cm,height=4.6cm,
  xlabel={quarter},ylabel={equity cash (USD m)},
  tick label style={font=\small},label style={font=\small},
  xmin=0.5,xmax=19.5,ymin=0,ymax=2.5,grid=major,grid style={gray!25},
  legend columns=3,legend style={font=\footnotesize,at={(0.5,-0.32)},anchor=north,draw=none,fill=none,/tikz/every even column/.append style={column sep=8pt}},
  table/col sep=comma]
\addplot[omDef,very thick,const plot mark mid,restrict x to domain=1:19] table[x=quarter,y=eq0]{figdata/markets-2/25-loans-clos-and-securitisation/paths.csv};
\addplot[omThm,very thick,const plot mark mid,restrict x to domain=1:19] table[x=quarter,y=eq1]{figdata/markets-2/25-loans-clos-and-securitisation/paths.csv};
\addplot[omProp,very thick,const plot mark mid,restrict x to domain=1:19] table[x=quarter,y=eq2]{figdata/markets-2/25-loans-clos-and-securitisation/paths.csv};
\legend{2\% a year,4.8\%,7\%}
\end{axis}
\end{tikzpicture}

\omcaption{The chapter's CLO under constant annual default rates of 2\%, 4.8\% and 7\%
(recovery 70\%, recoveries reinvested). Top: the overcollateralisation ratio at the BB
level against its trigger; once it fails, diversions hold it near the trigger for as
long as the interest allows. Bottom: the equity's quarterly cash, which the diversions
cut, to zero from quarter 12 at 7\% and to almost nothing from quarter 17 at 4.8\%; the last quarter's sale
of the loans is not shown. Illustrative; data: the chapter's
tutorial.}\label{fig:m2:loans-clos-and-securitisation:paths}
\end{omfigure}

\section{Who holds which tranche}

The notes are sold to investors whose needs match their risk: American CLO investors
are typically banks, mutual funds, insurance companies, pension funds and hedge funds.
The AAA notes pay a small spread over the floating rate for protection by 38\% of
subordination in the chapter's example, which suits investors that need high ratings
and floating-rate assets; the mezzanine notes pay more for thinner protection; the
equity is a leveraged position in the loans, about ten times, and earns the excess
spread, which suits investors who can bear large losses in a bad cycle. The manager is
paid fees senior to all the notes, whatever happens to the equity.

The structure makes the equity's returns very sensitive to defaults
(\cref{fig:m2:loans-clos-and-securitisation:irr}). Its internal rate of return falls
from 19.5\% without defaults to 13.1\% at 2\% a year and 5.0\% at 4\%, and it is
negative from about 5\%. The AAA notes, protected by the subordination and by the
tests, which turn the equity's cash into repayments, are paid in full in every case
the tutorial runs, even at 10\% a year.

\begin{omfigure}
\begin{tikzpicture}
\begin{axis}[width=11cm,height=5cm,
  xlabel={constant default rate (\% a year)},ylabel={equity IRR (\%)},
  tick label style={font=\small},label style={font=\small},
  xmin=0,xmax=10,ymin=-40,ymax=25,grid=major,grid style={gray!25},
  axis y line*=left,
  legend columns=2,legend style={font=\footnotesize,at={(0.5,-0.3)},anchor=north,draw=none,fill=none,/tikz/every even column/.append style={column sep=8pt}},
  table/col sep=comma]
\addplot[omDef,very thick,mark=*,mark size=1.2pt] table[x=cdr,y=irr]{figdata/markets-2/25-loans-clos-and-securitisation/irr.csv};
\addplot[gray,thin,forget plot] coordinates {(0,0) (10,0)};
\addplot[omThm,thick,dashed] coordinates {(4.78,-40) (4.78,25)};
\legend{equity IRR,equity first cut off (4.78\%)}
\end{axis}
\end{tikzpicture}

\omcaption{Internal rate of return of the chapter's CLO equity, over five years with
the loans sold at par at the end, against a constant annual default rate (recovery
70\%). The dashed line marks the default rate above which the tests cut the equity's
cash to zero in some quarter. Returns below $-40\%$, from 9\% a year, are off the
chart. Illustrative; data: the chapter's tutorial.}\label{fig:m2:loans-clos-and-securitisation:irr}
\end{omfigure}

\section{Tutorial: a CLO waterfall through a default wave}\label{tut:m2:loans-clos-and-securitisation:waterfall}

\begin{tutorial}
\textbf{Goal.} Build the waterfall, run the chapter's CLO through constant default
rates, watch the \omterm{def:m2:loans-clos-and-securitisation:tests}{overcollateralisation test} divert cash, and find the default rate at
which the equity stops being paid.
\textbf{End state:} \cref{fig:m2:loans-clos-and-securitisation:paths,fig:m2:loans-clos-and-securitisation:irr},
\cref{ex:m2:loans-clos-and-securitisation:base} and the numbers of the weekend problem.
\begin{tutsteps}
\item \textbf{The deal}: notes, triggers and the quarterly record.
\omcode{code/firm/waterfall/firm_waterfall.py}{15}{44}{Notes, the deal's terms and a period's record.}\label{lst:m2:loans-clos-and-securitisation:deal}
\item \textbf{The waterfall}: defaults, reinvestment, interest by class, tests and
diversions (\texttt{pay\_down} repays the notes in order), the final sale.
\omcode{code/firm/waterfall/firm_waterfall.py}{56}{94}{The quarterly waterfall.}\label{lst:m2:loans-clos-and-securitisation:run}
\item \textbf{The results}: the equity's internal rate of return and the lowest default
rate at which it is cut off, by bisection.
\omcode{code/firm/waterfall/firm_waterfall.py}{97}{117}{Equity IRR and the cut-off default rate.}\label{lst:m2:loans-clos-and-securitisation:irr}
\item \textbf{Run} \texttt{clo\_demo.base\_case()}, \texttt{clo\_demo.problem()} and
\texttt{fig\_clo.py}.
\end{tutsteps}
\textbf{What to change next.} Replace the constant default rate by a wave, 2\% a year
for two years, then 10\% for one, then 2\%, and compare the equity's return with the
constant rate that has the same cumulative defaults; then let the loans be sold at 95
instead of par at the end.
\end{tutorial}

\section{Build: the CLO waterfall engine}\label{bld:m2:loans-clos-and-securitisation:waterfall}

\begin{build}
\bfield{Purpose} The miniature firm's credit desk values CLO notes and equity under
default scenarios, and its risk team sees how close each deal is to its triggers.
\bfield{Interface} \texttt{Note(name, balance, spread, deferrable)};
\texttt{Deal(notes, par, loan\_spread, rate, fee, oc, ic, recovery, quarters,
reinvest\_quarters)}; \texttt{run(deal, cdr)} returning one \texttt{Period} per
quarter; \texttt{pay\_down(notes, amount)}; \texttt{equity\_irr(cash, invested)};
\texttt{cutoff\_cdr(deal)}.
\bfield{Rules} Quarterly; constant annual default rate; recoveries reinvested at par
during the reinvestment period, used to repay notes after it; fee senior to all notes;
deferrable junior interest; overcollateralisation cures limited to the interest
available; an interest-coverage failure diverts all remaining interest; loans sold at
par at the end.
\bfield{Acceptance tests} \texttt{code/firm/waterfall/tests/}: without defaults every
class is paid and the equity receives the excess spread and its principal; balances
and payments never go negative; above the cut-off default rate the equity is cut off in
some quarter, below it never.
\bfield{Stretch} Loan-level pools with ratings and a CCC bucket haircut in the tests;
default and prepayment scenarios from a credit model; the principal waterfall and
reinvestment criteria; valuing each class by simulation (One Quant Book 6).
\end{build}

\begin{omsources}
\item Federal Reserve Board, Financial Stability Report, May 2026.
\item M.~Guse, W.~Park, Z.~Saravay and Y.~Yook, ``Collateralized loan obligations in the
Financial Accounts of the United States'', FEDS Notes, September 2019.
\item NAIC Capital Markets Bureau, ``Middle market collateralized loan obligations
primer'', March 2025.
\item Chapman and Cutler, client alert on \emph{Loan Syndications and Trading
Association v.\ SEC} (D.C. Cir., 9 February 2018).
\end{omsources}

\section{Exercises}

\begin{exercise}[$\star$]\label{exo:m2:loans-clos-and-securitisation:1}
Compute the four overcollateralisation ratios of the chapter's CLO at issue.
\end{exercise}

\begin{exercise}[$\star$]\label{exo:m2:loans-clos-and-securitisation:2}
Why does a \omterm{def:m2:loans-clos-and-securitisation:loan}{leveraged loan}'s interest bill rise when the central bank raises rates, and
a fixed-rate bond's does not?
\end{exercise}

\begin{exercise}[$\star$]\label{exo:m2:loans-clos-and-securitisation:3}
What does a \omterm{def:m2:loans-clos-and-securitisation:loan}{covenant-lite} loan take away from its lenders?
\end{exercise}

\begin{exercise}[$\star\star$]\label{exo:m2:loans-clos-and-securitisation:4}
With a recovery of 70\%, how much of the portfolio must default before the A-level
test (trigger 1.18) fails?
\end{exercise}

\begin{exercise}[$\star\star$]\label{exo:m2:loans-clos-and-securitisation:5}
Why did the court find that the risk-retention rule did not fit open-market CLO
managers, and what incentive problem does risk retention address?
\end{exercise}

\begin{exercise}[$\star\star$]\label{exo:m2:loans-clos-and-securitisation:6}
The equity is about ten times levered on the loans. Using \cref{ex:m2:loans-clos-and-securitisation:base},
explain why its cash yield is 18.2\% when the loans pay 3.5\% over the floating rate.
\end{exercise}

\begin{exercise}[$\star\star\star$]\label{exo:m2:loans-clos-and-securitisation:7}
\emph{Coding.} End the reinvestment period after two years instead of five, so that
recoveries repay the notes, and find the new cut-off default rate. Explain the change.
\end{exercise}

\begin{exercise}[$\star\star\star$]\label{exo:m2:loans-clos-and-securitisation:8}
\emph{Find the flaw.} ``The AAA notes of a CLO are as safe as the AAA bonds of a
large company: same rating, same risk.'' Correct it.
\end{exercise}

\section{Problem: The Test That Trips}

\begin{problem}[{Weekend problem --- when the equity stops being paid}]\label{pb:m2:loans-clos-and-securitisation:1}
An investor considers the equity of the chapter's CLO: USD~500 million of loans at 3.5\%
over a floating rate of 4\%, notes of USD~450 million in five classes, USD~50 million of
equity, a senior fee of 0.45\%, overcollateralisation triggers of 1.25, 1.18, 1.10 and
1.05 after the AA, A, BBB and BB classes, an interest-coverage trigger of 1.20 after
the AA class, recoveries of 70\% reinvested for five years, and the loans sold at par
at the end.

\medskip\noindent\textbf{Part I --- The structure.}
\begin{enumerate}
\item What do the loans earn, and what do the fee and the notes cost, each year?
\item What does the equity receive each year without defaults, and what cash yield is
that?
\item What is the equity's return over five years without defaults?
\item How far can the par fall before the BB-level test fails?
\item What cumulative defaults is that, with a recovery of 70\%?
\end{enumerate}

\medskip\noindent\textbf{Part II --- Defaults.}
\begin{enumerate}[resume]
\item What is the equity's return at constant default rates of 2\%, 4\% and 6\% a year?
\item How much cash is diverted over five years at 4\%, 6\% and 10\%?
\item Find the lowest constant default rate at which the equity receives nothing in
some quarter.
\item At that rate, in which quarter is the equity first cut off, and how much par has
the portfolio lost?
\item What happens to the BB notes at 10\% a year?
\end{enumerate}

\medskip\noindent\textbf{Part III --- Levers.}
\begin{enumerate}[resume]
\item How does ending the reinvestment period after two years change the cut-off?
\item Why do the tests protect the AAA notes?
\item Why is a loan that recovers 70\% after default less damaging to the tests than a
bond that recovers 40\%?
\item What would a manager do, when the test is close to failing, to keep the equity
paid, and what are the risks?
\item Why does it matter for the equity that the loans are sold at par at the end?
\end{enumerate}

\medskip\noindent\textbf{Part IV --- Judgement.}
\begin{enumerate}[resume]
\item Which historical default rates would you compare 4.78\% with, and why is a
constant rate optimistic in a recession?
\item Who should hold the equity, and who the AAA notes?
\item How would you hedge the equity's exposure to a default wave?
\item State the \emph{named result}: the default rate at which the equity stops
receiving cash, and when.
\item In one sentence: what does the \omterm{def:m2:loans-clos-and-securitisation:tests}{overcollateralisation test} do?
\end{enumerate}
\end{problem}

\section{Interview questions}

\begin{interviewq}[$\star$ \iqroles{trader, researcher}]\label{iq:m2:loans-clos-and-securitisation:1}
What is a CLO, and who holds its tranches?
\end{interviewq}

\begin{interviewq}[$\star$ \iqroles{researcher, risk}]\label{iq:m2:loans-clos-and-securitisation:2}
How do overcollateralisation and \omterm{def:m2:loans-clos-and-securitisation:tests}{interest-coverage tests} work?
\end{interviewq}

\begin{interviewq}[$\star\star$ \iqroles{researcher}]\label{iq:m2:loans-clos-and-securitisation:3}
How would you value a CLO \omterm{def:m2:loans-clos-and-securitisation:tests}{equity tranche}?
\end{interviewq}

\begin{interviewq}[$\star\star$ \iqroles{trader, researcher}]\label{iq:m2:loans-clos-and-securitisation:4}
Compare a CLO tranche with a tranche of a \omterm{def:m2:credit-indices-and-tranches:index}{credit index}. What is different about its
risk?
\end{interviewq}

\begin{interviewq}[$\star\star$ \iqroles{risk, bank}]\label{iq:m2:loans-clos-and-securitisation:5}
What are the risks of the leveraged-loan market for financial stability?
\end{interviewq}

\begin{interviewq}[$\star\star\star$ \iqroles{developer}]\label{iq:m2:loans-clos-and-securitisation:6}
Design a system that runs the waterfalls of 2\,000 CLOs under 10\,000 scenarios each
night.
\end{interviewq}
